% app_e 附录G+H 共形变换+共形图 (书页 467-478 / p482-p493)
%--D-TAIL--
% 附录G Conformal Transformations
\chapter{共形变换}

% ---- 书页 467 / p482 ----
\textbf{共形变换}（conformal transformation）本质上是一种局部的标度改变。由于距离是靠度规来量度的，这类变换是通过用一个依赖于时空位置的（非零）函数去乘度规来实现的：
\begin{equation*}
\tilde{g}_{\mu\nu} = \omega^{2}(x) g_{\mu\nu},
\tag{G.1}
\end{equation*}
或等价地
\begin{equation*}
\widetilde{ds}^{2} = \omega^{2}(x) ds^{2},
\tag{G.2}
\end{equation*}
其中 $\omega(x)$ 是某个非零函数。（这里用 $x$ 表示时空坐标 $x^{\mu}$ 的全体。）注意，逆共形变换是平凡的：$g_{\mu\nu} = \omega^{-2}\tilde{g}_{\mu\nu}$。这类变换在广义相对论中有许多用途；我们最钟爱的两个用途是：在标量-张量理论中更换动力学变量（如 4.8 节所述），以及把时空重新映射成便于使用的共形图（见下一个附录）。

我们首先指出一个关键事实：\emph{类光曲线在共形变换下保持不变}。这句话的意思很简单：如果 $x^{\mu}(\lambda)$ 是一条关于 $g_{\mu\nu}$ 为类光的曲线，那么它关于 $\tilde{g}_{\mu\nu}$ 也是类光的。只要我们明白“曲线 $x^{\mu}(\lambda)$ 为类光，当且仅当它的切矢量 $\mathrm{d}x^{\mu}/\mathrm{d}\lambda$ 为类光”，这一点便立即得证：
\begin{equation*}
g_{\mu\nu}\frac{dx^{\mu}}{d\lambda}\frac{dx^{\nu}}{d\lambda} = 0.
\tag{G.3}
\end{equation*}
于是，在与它共形相关的度规中，我们有
\begin{equation*}
\tilde{g}_{\mu\nu}\frac{dx^{\mu}}{d\lambda}\frac{dx^{\nu}}{d\lambda}
= \omega^{2}(x) g_{\mu\nu}\frac{dx^{\mu}}{d\lambda}\frac{dx^{\nu}}{d\lambda} = 0.
\tag{G.4}
\end{equation*}
这样，按一个度规定义为类光的曲线，按任何与它共形相关的度规定义也是类光的。我们可以说，“共形变换保持光锥不变。”（事实上，你可以验证它们也保持任意两个四维矢量之间的夹角不变；我们的共形变换与复分析中人们熟悉的共形变换正共享这一特点。）

接下来让我们考虑几何量在共形变换下如何改变。共形变换不是坐标变换，而是几何本身的实际改变——例如，$\tilde{g}_{\mu\nu}$ 的类时测地线一般会不同于 $g_{\mu\nu}$ 的类时测地线。不过，我们可以利用共形变换来更换动力学变量：任何作为 $g_{\mu\nu}$ 的函数的量，都同样可以看作 $\tilde{g}_{\mu\nu}$ 与 $\omega(x)$ 的函数。这时我们说这些量被表达在\textbf{共形框架}（conformal frame）之中。在本附录里，我们汇集一些表达式，说明原度规 $g_{\mu\nu}$ 中的各量如何与共形度规 $\tilde{g}_{\mu\nu}$ 中的各量相联系。

% ---- 书页 468 / p483 ----
我们先从 Christoffel 符号开始考察。由于联络系数对度规的导数是线性的，对逆度规也是线性的，经共形变换后的联络具有如下形式：
\begin{equation*}
\widetilde{\Gamma}^{\rho}{}_{\mu\nu} = \Gamma^{\rho}{}_{\mu\nu} + C^{\rho}{}_{\mu\nu}.
\tag{G.5}
\end{equation*}
$C^{\rho}{}_{\mu\nu}$ 显然是一个张量，因为它是两个联络之差。显式计算表明它由下式给出：
\begin{equation*}
C^{\rho}{}_{\mu\nu} = \omega^{-1}\left(\delta^{\rho}_{\mu}\nabla_{\nu}\omega + \delta^{\rho}_{\nu}\nabla_{\mu}\omega - g_{\mu\nu}g^{\rho\lambda}\nabla_{\lambda}\omega\right).
\tag{G.6}
\end{equation*}
当我们考虑 Riemann 张量在共形变换下的行为时，这个公式立刻就派上了用场。事实上，对任何形如式 (G.5) 的联络改变，我们都有
\begin{equation*}
\widetilde{R}^{\rho}{}_{\sigma\mu\nu} = R^{\rho}{}_{\sigma\mu\nu} + \nabla_{\mu}C^{\rho}{}_{\nu\sigma} - \nabla_{\nu}C^{\rho}{}_{\mu\sigma} + C^{\rho}{}_{\mu\lambda}C^{\lambda}{}_{\nu\sigma} - C^{\rho}{}_{\nu\lambda}C^{\lambda}{}_{\mu\sigma}.
\tag{G.7}
\end{equation*}
于是，只需代入并埋头演算，就得到
\begin{align*}
\widetilde{R}^{\rho}{}_{\sigma\mu\nu} = R^{\rho}{}_{\sigma\mu\nu} &- 2\left(\delta^{\rho}_{[\mu}\delta^{\alpha}_{\nu]}\delta^{\beta}_{\sigma} - g_{\sigma[\mu}\delta^{\alpha}_{\nu]}g^{\rho\beta}\right)\omega^{-1}(\nabla_{\alpha}\nabla_{\beta}\omega) \\
&+ 2\left(2\delta^{\rho}_{[\mu}\delta^{\alpha}_{\nu]}\delta^{\beta}_{\sigma} - 2g_{\sigma[\mu}\delta^{\alpha}_{\nu]}g^{\rho\beta} + g_{\sigma[\mu}\delta^{\rho}_{\nu]}g^{\alpha\beta}\right)\omega^{-2}(\nabla_{\alpha}\omega)(\nabla_{\beta}\omega).
\tag{G.8}
\end{align*}
把第一个和第三个指标缩并，就得到 Ricci 张量，
\begin{align*}
\widetilde{R}_{\sigma\nu} = R_{\sigma\nu} &- \left[(n-2)\delta^{\alpha}_{\sigma}\delta^{\beta}_{\nu} + g_{\sigma\nu}g^{\alpha\beta}\right]\omega^{-1}(\nabla_{\alpha}\nabla_{\beta}\omega) \\
&+ \left[2(n-2)\delta^{\alpha}_{\sigma}\delta^{\beta}_{\nu} - (n-3)g_{\sigma\nu}g^{\alpha\beta}\right]\omega^{-2}(\nabla_{\alpha}\omega)(\nabla_{\beta}\omega),
\tag{G.9}
\end{align*}
其中 $n$ 是维数。提升一个指标（利用 $\tilde{g}^{\mu\nu} = \omega^{-2}g^{\mu\nu}$）并再作一次缩并，就得到标量曲率，
\begin{equation*}
\widetilde{R} = \omega^{-2}R - 2(n-1)g^{\alpha\beta}\omega^{-3}(\nabla_{\alpha}\nabla_{\beta}\omega) - (n-1)(n-4)g^{\alpha\beta}\omega^{-4}(\nabla_{\alpha}\omega)(\nabla_{\beta}\omega).
\tag{G.10}
\end{equation*}
另一个有用的量是标量场 $\phi$ 的协变导数。一阶协变导数在原框架与共形框架中是相等的，因为二者都等于偏导数：
\begin{equation*}
\widetilde{\nabla}_{\mu}\phi = \nabla_{\mu}\phi = \partial_{\mu}\phi.
\tag{G.11}
\end{equation*}

% ---- 书页 469 / p484 ----
不过，二阶导数涉及 Christoffel 符号，因此具有非平凡的变换律：
\begin{equation*}
\widetilde{\nabla}_{\mu}\widetilde{\nabla}_{\nu}\phi = \nabla_{\mu}\nabla_{\nu}\phi - \left(\delta^{\alpha}_{\mu}\delta^{\beta}_{\nu} + \delta^{\beta}_{\mu}\delta^{\alpha}_{\nu} - g_{\mu\nu}g^{\alpha\beta}\right)\omega^{-1}(\nabla_{\alpha}\omega)(\nabla_{\beta}\phi).
\tag{G.12}
\end{equation*}
我们可以用 $\tilde{g}^{\mu\nu}$ 对它缩并，得到 d'Alembert 算符，
\begin{equation*}
\widetilde{\Box}\phi = \omega^{-2}\Box\phi + (n-2)g^{\alpha\beta}\omega^{-3}(\nabla_{\alpha}\omega)(\nabla_{\beta}\phi).
\tag{G.13}
\end{equation*}
最后，我们也许想反其道而行，用共形度规来表达原度规中的各个量。这不过是件繁琐的计算活儿；为方便起见，这里把答案抄录如下。曲率张量及其缩并为
\begin{align*}
R^{\rho}{}_{\sigma\mu\nu} = \widetilde{R}^{\rho}{}_{\sigma\mu\nu} &+ 2\left(\delta^{\rho}_{[\mu}\delta^{\alpha}_{\nu]}\delta^{\beta}_{\sigma} - \tilde{g}_{\sigma[\mu}\delta^{\alpha}_{\nu]}\tilde{g}^{\rho\beta}\right)\omega^{-1}(\widetilde{\nabla}_{\alpha}\widetilde{\nabla}_{\beta}\omega) \\
&+ 2\tilde{g}_{\sigma[\mu}\delta^{\rho}_{\nu]}\tilde{g}^{\alpha\beta}\,\omega^{-2}(\widetilde{\nabla}_{\alpha}\omega)(\widetilde{\nabla}_{\beta}\omega),
\tag{G.14}
\end{align*}
\begin{align*}
R_{\sigma\nu} = \widetilde{R}_{\sigma\nu} &+ \left[(n-2)\delta^{\alpha}_{\sigma}\delta^{\beta}_{\nu} + \tilde{g}_{\sigma\nu}\tilde{g}^{\alpha\beta}\right]\omega^{-1}(\widetilde{\nabla}_{\alpha}\widetilde{\nabla}_{\beta}\omega) \\
&- (n-1)\tilde{g}_{\sigma\nu}\tilde{g}^{\alpha\beta}\,\omega^{-2}(\widetilde{\nabla}_{\alpha}\omega)(\widetilde{\nabla}_{\beta}\omega),
\tag{G.15}
\end{align*}
以及
\begin{equation*}
R = \omega^{2}\widetilde{R} + 2(n-1)\tilde{g}^{\alpha\beta}\omega(\widetilde{\nabla}_{\alpha}\widetilde{\nabla}_{\beta}\omega) - n(n-1)\tilde{g}^{\alpha\beta}(\widetilde{\nabla}_{\alpha}\omega)(\widetilde{\nabla}_{\beta}\omega),
\tag{G.16}
\end{equation*}
而标量场的协变导数则由下式给出：
\begin{equation*}
\nabla_{\mu}\nabla_{\nu}\phi = \widetilde{\nabla}_{\mu}\widetilde{\nabla}_{\nu}\phi + \left(\delta^{\alpha}_{\mu}\delta^{\beta}_{\nu} + \delta^{\beta}_{\mu}\delta^{\alpha}_{\nu} - \tilde{g}_{\mu\nu}\tilde{g}^{\alpha\beta}\right)\omega^{-1}(\widetilde{\nabla}_{\alpha}\omega)(\widetilde{\nabla}_{\beta}\phi)
\tag{G.17}
\end{equation*}
以及
\begin{equation*}
\Box\phi = \omega^{2}\widetilde{\Box}\phi - (n-2)\tilde{g}^{\alpha\beta}\omega(\widetilde{\nabla}_{\alpha}\omega)(\widetilde{\nabla}_{\beta}\phi).
\tag{G.18}
\end{equation*}

\section{习题}

% 习题编号照原书
\textbf{1.}\quad 证明共形变换保持类光测地线不变，也就是说，$g_{\mu\nu}$ 的类光测地线与 $\omega^{2}g_{\mu\nu}$ 的类光测地线相同。（我们已经知道它们保持类光\emph{曲线}不变；你还需要证明变换后的曲线仍然是测地线。）原度规与共形度规中的仿射参量之间是什么关系？

\textbf{2.}\quad 证明在二维情形下，总能找到一个共形变换（只要算符 $\nabla^{\mu}\nabla_{\mu}$ 可逆），使得变换后的度规的曲率至少在某个坐标卡片中为零。（一般无法在整个流形上同时做到这一点。）这意味着任何二维度规都可以局部地写成一个平直度规乘以一个共形因子。

\textbf{3.}\quad 假设两个度规由如下形式的整体共形变换相联系：
\begin{equation*}
\tilde{g}_{\mu\nu} = e^{\alpha(x)} g_{\mu\nu}.
\tag{G.19}
\end{equation*}

% ---- 书页 470 / p485 ----
\textbf{(a)}\quad 证明：如果 $\xi^{\mu}$ 是度规 $g_{\mu\nu}$ 的一个 Killing 矢量，那么它是度规 $\tilde{g}_{\mu\nu}$ 的一个共形 Killing 矢量。\textbf{共形 Killing 矢量}（conformal Killing vector）满足方程
\begin{equation*}
\nabla_{\mu}\xi_{\nu} + \nabla_{\nu}\xi_{\mu} = (\nabla_{\lambda}\alpha)\xi^{\lambda} g_{\mu\nu}.
\tag{G.20}
\end{equation*}

\textbf{(b)}\quad 证明 $\xi_{\mu}k^{\mu}$ 沿 $\tilde{g}_{\mu\nu}$ 中的光子测地线为常数。这里 $k^{\mu}$ 是光子的四维动量。

\textbf{(c)}\quad 证明共形时间 $\eta = \int dt/R(t)$ 与一个共形 Killing 矢量 $\xi = \partial_{\eta}$ 相联系。

\textbf{(d)}\quad 利用（c）小题重新导出尺度因子与红移之间的关系。

% 附录H Conformal Diagrams
\chapter{共形图}

% ---- 书页 471 / p486 ----
弯曲时空流形原则上可以复杂得超乎想象；幸运的是，我们常常可以用具有高度对称性（尤其是球对称性）的流形来近似物理上现实的情形。然而，即便是对称的时空，如果我们试图想象这类流形的整体结构，它们也会对我们的直观想象能力构成可怕的挑战。因此，能够画出一种标准化的时空图，用来刻画足够对称的时空的整体性质与因果结构，是很有用的。（所谓“因果结构”，指的是由光锥定义的不同事件的过去与未来之间的关系。）\textbf{共形图}（conformal diagrams，亦称 Carter--Penrose 图，或干脆叫彭罗斯图）优雅地实现了这一愿望。

共形图不过是一张普通的时空图，只不过我们对其度规实施了一个特别聪明的坐标变换。由于我们的目标是描绘时空的因果结构，而因果结构由光锥定义，所谓“聪明”，就是指新坐标 $x^{\mu'}$ 具有“类时”坐标与“径向”坐标，且径向光锥在时空图上可以始终如一地画成 $45^{\circ}$。此外，我们追求的是使“无限远”只离有限坐标值之遥的坐标，好让整个时空的结构一目了然。

正如前一个附录中所解释的，共形变换保持光锥不变。既然我们想找到光锥成 $45^{\circ}$ 的坐标，我们只需找到这样一组坐标，使得所关心的度规与另一个我们已知其光锥成 $45^{\circ}$ 的度规共形相关。（当然，光锥画成什么角度取决于我们所用的单位，或者说取决于我们如何画坐标轴；我们真正的意思，是取一组坐标 $T$、$R$，使得径向类光线满足 $dT/dR = \pm 1$。）

让我们从 Minkowski 空间开始，看看这个技巧如何施展。极坐标下的 Minkowski 度规为
\begin{equation*}
ds^{2} = -\mathrm{d}t^{2} + \mathrm{d}r^{2} + r^{2}d\Omega^{2},
\tag{H.1}
\end{equation*}
其中 $d\Omega^{2} = \mathrm{d}\theta^{2} + \sin^{2}\theta\,\mathrm{d}\phi^{2}$ 是单位二维球面上的度规。在这里，我们其实已经可以在每一处都把光锥画成 $45^{\circ}$（轨迹 $t = \pm r$ 是类光的），但我们希望通过换用取值范围有限的坐标，让整个时空的因果结构更加清晰。$\theta$、$\phi$ 坐标不会有什么异样，但我们会想仔细记下另外两个坐标的取值范围。

% ---- 书页 472 / p487 ----
首先当然有
\begin{equation*}
-\infty < t < \infty, \qquad 0 \le r < \infty.
\tag{H.2}
\end{equation*}
严格说来，世界线 $r = 0$ 代表一个坐标奇点，本应由另一张片（patch）来覆盖；但大家都明白这是怎么回事，所以我们就索性装作 $r = 0$ 表现良好。

第一种尝试（结果发现行不通）也许就是简单地重新缩放类时坐标与径向坐标，让它们覆盖一个有限的范围。一个不错的候选是利用反正切函数（如图 H.1 所示），定义 $\bar{t} = \arctan t$、$\bar{r} = \arctan r$。这样度规将取如下形式［利用 $\mathrm{d}\tan x = (1/\cos^{2}x)\mathrm{d}x$］：
\begin{equation*}
ds^{2} = -\frac{1}{\cos^{4}\bar{t}}\,\mathrm{d}\bar{t}^{2} + \frac{1}{\cos^{4}\bar{r}}\,\mathrm{d}\bar{r}^{2} + \tan^{2}\bar{r}\,d\Omega^{2},
\tag{H.3}
\end{equation*}
其取值范围为
\begin{equation*}
\begin{gathered}
-\frac{\pi}{2} < \bar{t} < \frac{\pi}{2} \\
0 \le \bar{r} < \frac{\pi}{2}.
\end{gathered}
\tag{H.4}
\end{equation*}
好消息是新坐标的取值范围有限；坏消息是光锥的斜率（由 $\mathrm{d}\bar{t}/\mathrm{d}\bar{r} = \pm\cos^{2}\bar{t}/\cos^{2}\bar{r}$ 给出）并不像我们希望的那样等于 $\pm 1$。如果我们据此画出相应的时空图（你不妨画着玩玩），类光线走的是哪条路将无从辨认，在时空的边缘处尤其如此。

走出这条死胡同的办法是：不再直来直去地摆弄原坐标 $t$ 和 $r$，而是变得更聪明一点，改用类光坐标：
\begin{equation*}
\begin{aligned}
u &= t - r \\
v &= t + r,
\end{aligned}
\tag{H.5}
\end{equation*}

%FIG{H.1}: {反正切函数把实线映射到一个有限区间。}

% ---- 书页 473 / p488 ----
%FIG{H.2}: {Minkowski 空间上的类光径向坐标。}

其相应的取值范围为
\begin{equation*}
-\infty < u < \infty, \qquad -\infty < v < \infty, \qquad u \le v.
\tag{H.6}
\end{equation*}
这些坐标如图 H.2 所示，图中每个点代表一个半径为 $r = \frac{1}{2}(v-u)$ 的二维球面。类光坐标下的 Minkowski 度规为
\begin{equation*}
ds^{2} = -\frac{1}{2}(\mathrm{d}u\,\mathrm{d}v + \mathrm{d}v\,\mathrm{d}u) + \frac{1}{4}(v-u)^{2}d\Omega^{2}.
\tag{H.7}
\end{equation*}
现在我们利用反正切把无限远拉到有限的坐标值处，令
\begin{equation*}
\begin{aligned}
U &= \arctan u \\
V &= \arctan v,
\end{aligned}
\tag{H.8}
\end{equation*}
其取值范围为
\begin{equation*}
-\pi/2 < U < \pi/2, \qquad -\pi/2 < V < \pi/2, \qquad U \le V.
\tag{H.9}
\end{equation*}
于是我们有
\begin{equation*}
\mathrm{d}u\,\mathrm{d}v + \mathrm{d}v\,\mathrm{d}u = \frac{1}{\cos^{2}U\cos^{2}V}(\mathrm{d}U\,\mathrm{d}V + \mathrm{d}V\,\mathrm{d}U),
\tag{H.10}
\end{equation*}
以及
\begin{align*}
(v-u)^{2} &= (\tan V - \tan U)^{2} = \frac{1}{\cos^{2}U\cos^{2}V}\left(\sin V\cos U - \cos V\sin U\right)^{2} \\
&= \frac{1}{\cos^{2}U\cos^{2}V}\sin^{2}(V-U),
\tag{H.11}
\end{align*}

% ---- 书页 474 / p489 ----
于是度规式 (H.7) 在这些坐标下成为
\begin{equation*}
ds^{2} = \frac{1}{4\cos^{2}U\cos^{2}V}\left[-2(\mathrm{d}U\,\mathrm{d}V + \mathrm{d}V\,\mathrm{d}U) + \sin^{2}(V-U)\,d\Omega^{2}\right].
\tag{H.12}
\end{equation*}
这种形式有一定的魅力，因为度规看起来就像一个相当简单的表达式乘以一个整体因子。我们还可以更进一步，通过如下变换回到一个类时坐标 $T$ 与一个径向坐标 $R$：
\begin{equation*}
T = V + U, \qquad R = V - U,
\tag{H.13}
\end{equation*}
其取值范围为
\begin{equation*}
0 \le R < \pi, \qquad |T| + R < \pi.
\tag{H.14}
\end{equation*}
现在度规为
\begin{equation*}
ds^{2} = \omega^{-2}(T,R)\left(-\mathrm{d}T^{2} + \mathrm{d}R^{2} + \sin^{2}R\,d\Omega^{2}\right),
\tag{H.15}
\end{equation*}
其中
\begin{align*}
\omega &= 2\cos U\cos V \\
&= 2\cos\left[\frac{1}{2}(T-R)\right]\cos\left[\frac{1}{2}(T+R)\right] \\
&= \cos T + \cos R.
\tag{H.16}
\end{align*}
因此，我们记作 $ds^{2}$ 的原 Minkowski 度规，可以看作经由一个共形变换与如下“非物理”度规相联系：
\begin{equation*}
\begin{aligned}
\widetilde{ds}^{2} &= \omega^{2}(T,R)\,ds^{2} \\
&= -\mathrm{d}T^{2} + \mathrm{d}R^{2} + \sin^{2}R\,d\Omega^{2}.
\end{aligned}
\tag{H.17}
\end{equation*}
这描述的是流形 $\mathbf{R} \times S^{3}$，其中的三维球面是纯类空的、完全浑圆的，且不随时间改变。这个度规是有曲率的，这与 Minkowski 时空不同。但这本不该困扰我们，因为它是非物理的；经由共形变换得到的真正物理的度规，无论我们选择什么坐标，都恰恰是平直时空。事实上，度规式 (H.17) 正是“Einstein 静态宇宙”的度规——它是 Einstein 方程带一个理想流体和一个宇宙学常数的静态解（图 H.3）。当然，$\mathbf{R} \times S^{3}$ 上坐标的完整取值范围通常是 $-\infty < T < \infty$、$0 \le R \le \pi$，而 Minkowski 空间则被映射到由式 (H.14) 定义的子空间中。整个 $\mathbf{R} \times S^{3}$ 可以画成一个柱面，其中每个 $T$ 为常数的圆周代表一个三维球面。阴影区域代表 Minkowski 空间。我们可以把阴影区域摊开，把 Minkowski 空间表示成一个三角形，如图 H.4 所示。这就是共形图。图中每个点代表一个二维球面。

% ---- 书页 475 / p490 ----
%FIG{H.3}: {Einstein 静态宇宙 $\mathbf{R} \times S^{3}$，描绘成一个柱面。阴影区域与 Minkowski 空间共形相关。}

事实上，Minkowski 空间只是上图（包括 $R = 0$ 在内）的\emph{内部}；边界并不属于原来的时空。这些边界被称为\textbf{共形无限远}（conformal infinity），而原时空与共形无限远之并集则称为\textbf{共形紧化}（conformal compactification），它是一个带边流形。共形图的结构使我们可以把共形无限远再细分成几个不同的区域：

%FIG{H.4}: {Minkowski 空间的共形图。整幅图中光锥都张成 $\pm 45^{\circ}$。}

% ---- 书页 476 / p491 ----
\begin{equation*}
\begin{aligned}
i^{+} &= \text{未来类时无限远}\ (T=\pi,\ R=0)\\
i^{0} &= \text{类空无限远}\ (T=0,\ R=\pi)\\
i^{-} &= \text{过去类时无限远}\ (T=-\pi,\ R=0)\\
\mathscr{I}^{+} &= \text{未来类光无限远}\ (T=\pi-R,\ 0<R<\pi)\\
\mathscr{I}^{-} &= \text{过去类光无限远}\ (T=-\pi+R,\ 0<R<\pi)
\end{aligned}
\end{equation*}
（$\mathscr{I}^{+}$ 与 $\mathscr{I}^{-}$ 分别读作“scri-plus”与“scri-minus”。）注意，$i^{+}$、$i^{0}$ 与 $i^{-}$ 其实是\emph{点}，因为 $R = 0$ 与 $R = \pi$ 是 $S^{3}$ 的北极和南极。而 $\mathscr{I}^{+}$ 与 $\mathscr{I}^{-}$ 则实际上是类光曲面，具有 $\mathbf{R} \times S^{2}$ 的拓扑。

Minkowski 时空的共形图包含若干重要特征。图中径向类光测地线成 $\pm 45^{\circ}$。所有类时测地线都始于 $i^{-}$、终于 $i^{+}$；所有类光测地线都始于 $\mathscr{I}^{-}$、终于 $\mathscr{I}^{+}$；所有类空测地线则既始于 $i^{0}$ 又终于 $i^{0}$。另一方面，也可以存在终于类光无限远的非测地类时曲线，只要它们变得“渐近类光”。

能把整个 Minkowski 空间塞进一小张纸片上固然惬意，但其实我们并没有由此学到多少原来不知道的东西。共形图真正更有用武之地，是当我们想表示稍微复杂一些的时空的时候，比如黑洞的时空。如第 6 章所讨论的，渐近平直时空（或时空的一个区域）是指那些其 $\mathscr{I}^{+}$、$i^{0}$ 与 $\mathscr{I}^{-}$ 的结构与 Minkowski 空间相同的时空。同样重要的是，共形图让我们得以了解时空的因果结构，比如两个指定点的过去或未来光锥是否相交。在 Minkowski 空间中，任何两点都是如此；但弯曲时空可以更有趣，正如我们在第 2 章膨胀宇宙的例子中看到的那样。

让我们考虑第 2 章引入的那个宇宙学时空的共形图，它生动地展示了这一技术的用处。在空间上取极坐标后，度规变为
\begin{equation*}
ds^{2} = -\mathrm{d}t^{2} + t^{2q}\left(\mathrm{d}r^{2} + r^{2}d\Omega^{2}\right),
\tag{H.18}
\end{equation*}
其中我们对尺度因子选择了幂律行为 $a(t) = t^{q}$，且 $0 < q < 1$。这个度规与 Minkowski 空间度规的一个关键差别，在于 $t = 0$ 处存在奇点，它限制了坐标的取值范围：
\begin{equation*}
0 < t < \infty
\tag{H.19}
\end{equation*}
\begin{equation*}
0 \le r < \infty.
\tag{H.20}
\end{equation*}
除了这一受限制的坐标范围之外，我们的分析几乎完全照搬平直时空的情形。这是因为我们可以把度规式 (H.18) 化成平直时空乘以一个共形因子的形式；一旦做到这一点，我们只需重复之前的坐标变换，就能把膨胀宇宙的度规表达为 Einstein 静态宇宙的度规乘以一个共形因子。

% ---- 书页 477 / p492 ----
我们首先选取一个新的时间坐标 $\eta$，它有时被称为\textbf{共形时间}（conformal time），满足
\begin{equation*}
\mathrm{d}t^{2} = t^{2q}\mathrm{d}\eta^{2},
\tag{H.21}
\end{equation*}
亦即
\begin{equation*}
\eta = \frac{1}{1-q}\,t^{1-q}.
\tag{H.22}
\end{equation*}
这一简单选择让我们得以把尺度因子提取成一个整体共形因子：
\begin{equation*}
ds^{2} = \left[(1-q)\eta\right]^{2q/(1-q)}\left(-\mathrm{d}\eta^{2} + \mathrm{d}r^{2} + r^{2}d\Omega^{2}\right).
\tag{H.23}
\end{equation*}
$\eta$ 的取值范围与 $t$ 相同，
\begin{equation*}
0 < \eta < \infty.
\tag{H.24}
\end{equation*}
注意，$\eta$ 是一个类时坐标［意义是矢量 $\partial_{\eta}$ 是类时的，$ds^{2}(\partial_{\eta},\partial_{\eta}) < 0$］，但它并不量度共动时钟（空间坐标为常数的时钟）的固有时。如果我们考虑轨迹 $x^{\mu}(\lambda) = (\eta(\lambda), 0, 0, 0)$，并计算固有时 $\tau(\eta)$，就会发现它等于我们先前的那个时间坐标，而不是新的这个：$\tau \propto t \propto \eta^{1/(1-q)}$。所以 $\eta$ 是类时坐标，却不是任何人实际会测到的时间。这完全没有问题，恰好可以作为一个例证，说明可观测量与时空坐标这两个概念是彼此独立的。

既然膨胀宇宙的度规已经写成了共形因子乘 Minkowski 度规的形式，我们就可以施行同一串坐标变换——式 (H.5)、(H.8) 与 (H.13)——只不过让 $\eta$ 顶替 $t$ 的位置。这些变换把我们的坐标从 $(\eta, r)$ 变到 $(T, R)$，此时的取值范围为
\begin{equation*}
0 \le R, \qquad 0 < T, \qquad T + R < \pi.
\tag{H.25}
\end{equation*}
度规式 (H.23) 变为
\begin{equation*}
ds^{2} = \omega^{-2}(T,R)\left(-\mathrm{d}T^{2} + \mathrm{d}R^{2} + \sin^{2}R\,d\Omega^{2}\right),
\tag{H.26}
\end{equation*}
其中，经过对三角恒等式的一番英勇运用，可以发现共形因子形如
\begin{equation*}
\omega(T,R) = \left(\frac{\cos T + \cos R}{2\sin T}\right)^{2q}(\cos T + \cos R).
\tag{H.27}
\end{equation*}
共形因子的精确形状其实并非头等重要；关键的特征是，我们又一次把度规表达成了 Einstein 静态宇宙的度规乘以一个共形因子。这个情形与平直时空情形的重要区别在于，类时坐标终止于 $T = 0$ 处的奇点；除此之外，时空图完全相同。

% ---- 书页 478 / p493 ----
于是我们得到图 H.5 的共形图，它看上去像 Minkowski 图（图 H.4）的上半部分。光锥依旧张成 $45^{\circ}$。我们看到共形图如何让因果结构一目了然：很容易在时空中选出两个事件，使得它们的过去光锥在相交之前先撞上奇点（而未来光锥则总会交叠）。对于更复杂的几何，这种表示时空的便利方式将更加有用。

%FIG{H.5}: {Robertson--Walker 宇宙（$a(t) \propto t^{q}$，$0<q<1$）的共形图。虚线代表 $t = 0$ 处的奇点（它同时对应于 $T = 0$）。}
